\begin{theorem}[Simply connected rectangular cell regions]%
    \label{thm:peps_torus_rectangle_realization}
    \lean{TNLean.PEPS.isSimplyConnected_torusRegionRealization_rectangle}
    \leanok
    \uses{lem:peps_torus_region_realization,
        lem:peps_torus_rectangle_cut_connectivity}
    Let the torus periods be positive integers $w,h$. Let $a,b\geq0$ and
    $\ell,m>0$ be integers with $\ell<w$, $m<h$, $a+\ell\leq w$, and
    $b+m\leq h$, and let $R$ be the native rectangle with coordinate ranges
    $a\leq x<a+\ell$ and $b\leq y<b+m$.
    Then the actual closed-cell realization of $R$ is simply connected.
    This is auxiliary geometry for the stripe reduction in
    \cite[Theorem~6.7]{Schuch2010PEPS}, local source lines 1995--2011.
\end{theorem}

\begin{proof}\leanok
    The torus projection maps the planar rectangle
    $[a-1/2,a+\ell-1/2]\times[b-1/2,b+m-1/2]$ onto the closed-cell
    realization of $R$. To verify surjectivity, choose in each coordinate
    an integer center in the indicated range at distance at most $1/2$;
    conversely, every cell centered in $R$ lies in this planar rectangle.
    Since its side lengths are smaller than the periods, projection is
    injective there. The continuous bijection from this compact rectangle
    to its Hausdorff image is a homeomorphism. The planar rectangle is
    nonempty and convex, hence contractible and simply connected.
\end{proof}
